\section{Bridges to other areas and to the author's other workbenches}\label{sec:overlaps}

In the author's view the connections of this research to other areas are among its most valuable outcomes. This section first treats the connections of the period system to established theories, and then the typed edges to the author's other workbenches. For each it states what is established and what is open.

\subsection{The period system and established theories}\label{ssec:theories}

\paragraph{The paramodular group of level $6$.} Gritsenko and Hulek define, for the alternating form $W_t$ of type $(1,t)$ on $\ZZ^4$, the integral paramodular group $\widetilde\Gamma_t=Sp(W_t,\ZZ)$, and they realise it as a group $\Gamma_t$ of rational symplectic matrices \cite[§1]{GritsenkoHulek}. Two alternating forms on $\ZZ^4$ with the same elementary divisors are equivalent under $GL_4(\ZZ)$ (the Frobenius normal form of alternating integral matrices), and $Q_0$ has elementary divisors $(1,1,6,6)$ (Section~\ref{sec:monodromy}). So $Sp(Q_0,\ZZ)\cong\widetilde\Gamma_6$: the ambient group of the period system is the paramodular group of level $6$. It contains $G_{\ZZ}$ as the stabiliser of the primitive isotropic vector $\gamma$, whose divisor $Q_0(\gamma,V)$ is $\ZZ$, and $G_{\ZZ}$ contains the Jacobi group $SL_2(\ZZ)\ltimes H(\ZZ)$ with index $6$ (Section~\ref{sec:monodromy}).

The character $e(\Phi/12)$, with $e(x)=e^{2\pi ix}$, has the following restrictions:
\begin{itemize}
\item on the Levi group, $e(\operatorname{ab}(M)/12)=v_\eta^2(M)$, where $v_\eta$ is the multiplier of Dedekind's $\eta$. Indeed $v_\eta^2$ is the multiplier of the weight-one form $\eta^2$, hence a character of $SL_2(\ZZ)$, and $v_\eta^2(T)=e(1/12)$ because $\eta(\tau+1)=e(1/24)\eta(\tau)$; two characters of $SL_2(\ZZ)$ that agree on $T$ are equal, since $T$ generates the abelianisation;
\item on the standard Heisenberg group $H(\ZZ)=\{h(p,r,6\kappa)\}$, identified with the integral Heisenberg group by $h(p,r,6\kappa)\mapsto[p,r;\kappa]$, it is $v_H[p,r;\kappa]=(-1)^{p+r+pr+\kappa}$ (Proposition~\ref{prop:heis}(d)), which is the Heisenberg multiplier of the odd Jacobi theta function in \cite[(1.4)--(1.5)]{GritsenkoHulek};
\item on the centre, $e(\Phi(h(0,0,s))/12)=e(s/12)$.
\end{itemize}
For $t=6$, Lemma~1.1 of \cite{GritsenkoHulek} gives characters $\chi_3$ and $\chi_4$, of orders $3$ and $4$, of the extended group $\Gamma_6^+$. On the Jacobi parabolic of their realization they restrict to $v_\eta^{8}$ and $v_\eta^{6}$ on $SL_2(\ZZ)$, to $v_H^{4}=1$ and $v_H^{3}=v_H$ on $H(\ZZ)$, and to $e(\kappa/3)$ and $e(\kappa/4)$ on the central elements $[0,0;\kappa/6]$. Hence $\chi_{12}=\chi_3\chi_4^{-1}$ has order $12$ and restricts there to $v_\eta^2$, $v_H$ and $e(\kappa/12)$: the same three restrictions as $e(\Phi/12)$. The Lemma was read through an extraction of its statement, not in the full text. Whether the parabolic of \cite{GritsenkoHulek} is conjugate in $\Gamma_6^+$ to $G_{\ZZ}$, and whether $e(\Phi/12)$ is the restriction of $\chi_{12}$ to $G_{\ZZ}$, is not verified here (Question~\ref{q:paramodular}). My assessment: I think the identification is likely, because the two characters have the same order and the same restrictions to the three pieces; what has to be computed is the value of $\chi_{12}$ on the two minimal positive transvections $L_T$ (along $u$, of divisor $6$) and $t_\gamma$ (along $\gamma$, of divisor $1$), which both have $\Phi=1$. Since $G_{\ZZ}^{\mathrm{ab}}$ is generated by the classes of these two elements (Theorem~\ref{thm:mono}(f)), if the two values are equal and of order $12$, then $G=G_{\ZZ}\cap\ker\chi_{12}$, and the monodromy group of the $S^6$ period system is cut out of the stabiliser of $\gamma$ in the paramodular group by a character of Gritsenko and Hulek.

\paragraph{Spin structures and theta characteristics.} Quadratic refinements of the mod-$2$ intersection form of a surface are its spin structures, with the Arf invariant as their parity \cite{Johnson}; on a Riemann surface they are its theta characteristics \cite{AtiyahSpin}. On a torus the odd spin structure is the one fixed by the mapping class group $SL_2(\ZZ)$, and Proposition~\ref{prop:heis}(b) shows that $SL_2$-invariance forces it in the monodromy group. The record's higher-torus update exhibits a Lie-framed torus whose spin form $x+xy+y$ has Arf invariant one, and this is the odd refinement (Section~\ref{sec:monodromy}). Whether that torus's first homology is the lattice $\langle\bar u,\bar w\rangle$ is Question~\ref{q:arfone}. My assessment: if it is, the update realises geometrically the spin structure that cuts out the Heisenberg part of the monodromy, and the update's construction and the monodromy computation describe the same object.

\paragraph{The cusp.} The cusp monodromy is $T_0=I+N$ with $N^2=0$ and $\operatorname{rank}N=2$. On $\Lambda^2V$ the induced nilpotent $N_2(a\wedge b)=Na\wedge b+a\wedge Nb$ satisfies $N_2^2(a\wedge b)=2\,Na\wedge Nb$ and $N_2^3=0$, and $N_2^2(w\wedge\delta)=-2\,u\wedge\gamma\neq0$. So $N_2^2\neq0=N_2^3$. For $H^2$ of a semistable degeneration of surfaces with trivial canonical class, $N^2\neq0$ is the monodromy of Kulikov's type~III \cite{Kulikov,PerssonPinkham}. The manuscript describes its central fibre at the cusp as a degree-six del Pezzo surface glued to itself along opposite sides of its hexagon, with two triple points (Section~\ref{sec:candidate}); its dual complex then has one vertex, three edges and two triangles, a triangulation of the torus (Euler characteristic $1-3+2=0$). This is the combinatorial type of Mumford's degenerations of abelian surfaces. The dictionary of Kulikov and of Persson and Pinkham rests on polarised limits, and $(V,Q_0,N)$ supports none, since $b_N$ is indefinite (Section~\ref{ssec:mononeg}). So the comparison holds for the Jordan type and the combinatorics. My assessment: I think the combinatorial match indicates that the degeneration is of Mumford type, and a substitute for the polarisation is what a Hodge-theoretic comparison needs.

\paragraph{The modular part.} Up to the factor $6$ in the induced form, the Levi group acts on $\gamma^\perp/\langle\gamma\rangle$ as $SL_2(\ZZ)$ acts on the first homology of an elliptic curve, and the Heisenberg part records the extension data. Theorem~\ref{thm:mono} and Proposition~\ref{prop:heis} describe both exactly. Whether $\gamma^\perp/\langle\gamma\rangle$ is the homology of an elliptic family inside or under the tori is Question~\ref{q:elliptic}.

\paragraph{Arithmetic, not thin.} $G$ has finite index in the integral points of its Zariski closure (Theorem~\ref{thm:mono}(h)), so it is not thin in the sense of Sarnak \cite{SarnakThin}; Brav and Thomas exhibit hypergeometric monodromy groups in $Sp(4,\ZZ)$ of the thin kind \cite{BravThomas}. Arithmeticity makes the methods for congruence subgroups available, and they give the orbit classification of Theorem~\ref{thm:orbitsS6}.

\subsection{The edges to the author's other workbenches}

The audit note \note{47\_} records every crossing between the owner's workbenches that the audit found, as typed edges: \emph{P}, a proof-level import; \emph{I}, the same identity or object on both sides; \emph{A}, an analogy or a shared number only. The table lists the edges that touch the $S^6$ material.

\begin{center}\small
\begin{longtable}{@{}>{\raggedright\arraybackslash}p{0.13\linewidth}>{\raggedright\arraybackslash}p{0.44\linewidth}>{\raggedright\arraybackslash}p{0.07\linewidth}>{\raggedright\arraybackslash}p{0.28\linewidth}@{}}
\toprule
Edge & What crosses & Type & Status\\
\midrule
\endhead
$S^6\to$ YM & The imaginary part of the period matrix, used as a magnetic background in the Yang--Mills workbench. Only $D_{\mathrm{YM}}=L_Tq_T+6m_T^2=-\det_{\R}\Pi$ enters (Section~\ref{sec:candidate}) & P & audited (\cite{YMPaper}, §§6.3 and~8)\\
$S^6\to$ ES & The integral monodromies $A_j=(T_j^{-1})^{\mathsf T}$ of the $(3,4,\infty)$ period system, in the Erdős--Straus package ES-09, with $A_1^3=A_2^4=A_1A_2A_\infty=I$. Nothing else from the construction is imported & P & the orbits at every level: Theorem~\ref{thm:orbitsS6}; the package's statements are correct for odd levels, which the application uses, and do not hold as stated at even levels \cite[§5.2]{ESPaper}\\
ES $\leftrightarrow$ $S^6$ & The Niemeier lattice with root system $A_5^4D_4$ and glue index $72$: the Erdős--Straus result R10, and S6-4 & I & the Erdős--Straus glue code verified; the $S^6$ code is in the record's archive \fn{28\_}, and the two were not compared\\
$S^6$ and the Jacobian map $\to$ NS & Fluid coordinates, and gauge-scalar equations from the $S^6$ material carried to Navier--Stokes on $\mathbb T^3$; a Beltrami-mode forced solution from the $S^6$ oscillator & P & the forced solution verified exactly \cite{NSPaper}; the rest located (audit note \note{47\_}, edge 13)\\
$S^6\leftrightarrow$ ES & The number $5184=72^2$. It is $\det I$ in S6-5, where $I$ has Gram matrix $6G_{D_4}$, so $\det I=6^4\det D_4$; and it is $|\operatorname{disc}(A_5^4D_4)|=(\det A_5)^4\det D_4$ in S6-4 and R10 & A & both products $6^4\cdot4$ checked; no map between the two lattices was found\\
\bottomrule
\end{longtable}
\end{center}

Two further remarks. The Yang--Mills repository has an \texttt{s6/} topic entry, which holds the key-advances reader, and the record's frozen edition of 9 September was released together with companion Yang--Mills and Navier--Stokes editions (guide \fn{29\_}). The $S^6$ construction and the Jacobian counterexample are different objects, as the attribution file of the Yang--Mills repository states; the Jacobian counterexample is treated in its own paper \cite{JacobianPaper}. No crossing was found between the correction note and the zeta programme: the derived criterion (Theorem~\ref{thm:derived}) and the local algebra of Section~\ref{sec:correction} use nothing from the other workbenches.
